\subsection{Математическое понятие центра масс}
%%%%%%%%%%%%%%%%%%%%%%%%%%%%%%%%%%%%%%%%%%%%%%%%%%%%
%%%%%%%%%%%%%%%%%%%%%%%%%%%%%%%%%%%%%%%%%%%%%%%%%%%%
%%%%%%%%%%%%%%%%%%%%%%%%%%%%%%%%%%%%%%%%%%%%%%%%%%%%	
 \z[{\begin{wrapfigure}[6]{r}{0.3\textwidth}%
		\vspace{-20pt}%
		\begin{flushright}%
 			\includegraphics[width=0.3\textwidth]{6.2.1}%
		\end{flushright}%
	\end{wrapfigure}%
	}]%
	{Однородная пластина имеет форму полукруга радиуса $r$, соединенного с прямоугольником длины $2r$ и ширины $h$ (см. рисунок). Найдите отношение $h/r$, если центр масс всей системы совпадает с геометрическим центром полукруга (точкой $C$). Расстояние от центра масс полукруга $C_1$ до его геометрического центра $C$ равно $\frac{4r}{3\pi}$.}
 %
{$\frac{h}{r} = \sqrt{\frac{2}{3}}.$}
%%%%%%%%%%%%%%%%%%%%%%%%%%%%%%%%%%%%%%%%%%%%%%%%%%%%	
 \z[{\begin{wrapfigure}[7]{r}{0.3\textwidth}%
		\vspace{-15pt}%
		\begin{flushright}%
 			\includegraphics[width=0.3\textwidth]{6.2.2}%
		\end{flushright}%
	\end{wrapfigure}%
	}]%
	{В однородном тонком диске радиуса $R$ вырезано отверстие радиуса $r < R/2$, центр которого находится на расстоянии $R/2$ от центра диска (см. рисунок). На каком расстоянии $x$ от центра диска находится центр масс этой системы?}
 %
{$x = \frac{r^2 R}{2(R^2 - r^2)}$}
%%%%%%%%%%%%%%%%%%%%%%%%%%%%%%%%%%%%%%%%%%%%%%%%%%%% 
 \z{Внутри однородного шара радиусом $R = 39$ см находится сферическая полость радиусом $r = 13$ см, касающаяся поверхности шара. На каком расстоянии от центра шара находится центр масс шара с полостью? Ответ выразить в сантиметрах.}
%
{$x = \frac{R - r}{\left(\frac{R}{r}\right)^3 - 1} = 1\,\text{см}$}
%%%%%%%%%%%%%%%%%%%%%%%%%%%%%%%%%%%%%%%%%%%%%%%%%%%% 
 \z{Определить центр масс треугольника, если масса равномерно распределена:
 	\pu в вершинах
 	\pu по сторонам
 	\pu по площади
	}
%
{точка пересечения: \puan медиан; \puan биссектрис сер. тр.; \puan медиан.}
%%%%%%%%%%%%%%%%%%%%%%%%%%%%%%%%%%%%%%%%%%%%%%%%%%%% 
 \z{Одинакового размера шарики массой $m$, $3m$, $5m$, $\dots$, $2019m$ расположены вдоль прямой. Где находится центр масс таких шариков?}
%
{$2691R \text{ от центра шарика массой } m$}